%=====================================================================
\chapter{Feature Engineering and Ensembles}\label{ch:features}
%=====================================================================
\locator{3}{Part III --- Machine Learning Core}

\begin{outcomes}
By the end of this chapter you will be able to:
\begin{tight}
  \item \textbf{Construct} informative features from raw columns using ratios, aggregations, time deltas and domain knowledge.
  \item \textbf{Apply} filter, wrapper and embedded methods of feature selection.
  \item \textbf{Explain} how bagging reduces variance and boosting reduces bias.
  \item \textbf{Distinguish} random forests from gradient boosting and choose between them.
  \item \textbf{Interpret} feature importances, and state the traps in doing so.
\end{tight}
\end{outcomes}

\begin{prereq}
\chref{classification} (decision trees --- the base learner for every ensemble
here) and \chref{evaluation} (you must be able to measure before you can improve).
\end{prereq}

\begin{keyterms}
feature engineering $\bullet$ feature selection $\bullet$ filter / wrapper /
embedded methods $\bullet$ curse of dimensionality $\bullet$ ensemble $\bullet$
bagging $\bullet$ bootstrap $\bullet$ random forest $\bullet$ boosting $\bullet$
gradient boosting $\bullet$ stacking $\bullet$ feature importance $\bullet$
permutation importance
\end{keyterms}

\section{Features Beat Algorithms}

\begin{conceptbox}[The Most Reliable Lesson in Applied Machine Learning]
Given a fixed dataset, moving from logistic regression to a tuned gradient
boosting machine might buy you two or three points of performance. Adding one good
feature that encodes something a domain expert knows can buy you twenty. Spend your
time accordingly.
\end{conceptbox}

\rowcolors{2}{white}{lightbg}
\begin{longtable}{@{}p{2.9cm}p{4.7cm}p{6.9cm}@{}}
\toprule
\rowcolor{primary!15}
\textbf{Technique} & \textbf{Construction} & \textbf{Example} \\
\midrule
\endhead
Ratios & Divide one column by another & \texttt{submissions / days\_enrolled} normalises for late joiners \\
Differences \& deltas & Subtract, or compare with a previous row & \texttt{this\_week\_logins $-$ last\_week\_logins} captures \emph{change}, which raw counts hide \\
Aggregations & Group and summarise & mean, max and standard deviation of a sensor over a 1-hour window \\
Time features & Extract components & hour of day, day of week, is\_weekend --- crucial in security and IoT \\
Counts of rare events & Count occurrences in a window & failed logins in the last 5 minutes \\
Binning & Convert numeric to ordinal & attendance into \{low, medium, high\} for a readable rule \\
Interactions & Multiply two features & \texttt{low\_attendance $\times$ first\_generation} \\
Domain flags & Encode expert rules & \texttt{missed\_two\_consecutive\_labs} \\
\bottomrule
\end{longtable}

\begin{worked}[Turning Raw Logs into Features]
You have one row per VLE click: \texttt{(student, timestamp, resource)}. A model
needs one row per student, using weeks 1--4 only.

\smallskip
\textbf{Volume:} \texttt{n\_clicks}, \texttt{n\_distinct\_resources},
\texttt{n\_active\_days}.

\textbf{Recency:} \texttt{days\_since\_last\_click}. Frequently the single
strongest predictor of disengagement --- and note it is unobtainable from any
aggregate count.

\textbf{Trend:} \texttt{clicks\_wk4 $-$ clicks\_wk1}, and the slope of a line
fitted through the four weekly totals. A student at 200 clicks and falling is in
more danger than one at 80 and steady, yet both have similar totals.

\textbf{Rhythm:} \texttt{share\_of\_clicks\_after\_midnight},
\texttt{n\_weekend\_days\_active}. Behavioural signatures, not effort measures.

\textbf{Consistency:} standard deviation of daily clicks. Steady beats bursty.

\smallskip
\textbf{From three raw columns, twelve features} --- and the ones most likely to
matter (recency, trend, consistency) do not appear in the raw data at all. This is
what feature engineering is. Everything here uses only weeks 1--4, so framing
question 4 (\chref{lifecycle}) is satisfied by construction.
\end{worked}

\section{Feature Selection}

\begin{alertbox}[The Curse of Dimensionality]
As the number of features grows, the volume of the feature space explodes and your
data becomes sparse within it --- every point ends up roughly equidistant from
every other, which destroys distance-based methods. More features is not more
information; past a point it is more noise.
\end{alertbox}

\dsfig{%
\begin{tikzpicture}[font=\scriptsize,
  m/.style={rounded corners=3pt, draw=#1, fill=#1!8, text width=3.9cm, align=center,
            font=\scriptsize, inner sep=5pt, minimum height=2.9cm},
  hd/.style={rounded corners=3pt, fill=#1, text=white, font=\scriptsize\bfseries,
             text width=3.9cm, align=center, minimum height=0.6cm}]

\node[hd=secondary] at (0,0)   {FILTER};
\node[m=secondary]  at (0,-1.85) {%
  Score each feature \emph{on its own}, before any model.\\[4pt]
  \emph{correlation, chi-square, mutual information}\\[4pt]
  {\color{greenhl!45!black}\textbf{+} very fast, model-agnostic}\\[2pt]
  {\color{accent}\textbf{--} blind to combinations: two useless-alone features can be
   decisive together}};

\node[hd=greenhl] at (4.5,0)   {WRAPPER};
\node[m=greenhl]  at (4.5,-1.85) {%
  Try subsets, train a model on each, keep the best.\\[4pt]
  \emph{forward selection, backward elimination, RFE}\\[4pt]
  {\color{greenhl!45!black}\textbf{+} accounts for interactions; tailored to your model}\\[2pt]
  {\color{accent}\textbf{--} expensive; overfits the selection set unless nested CV is used}};

\node[hd=purplehl] at (9.0,0)  {EMBEDDED};
\node[m=purplehl]  at (9.0,-1.85) {%
  Selection happens \emph{during} training.\\[4pt]
  \emph{lasso (\chref{regression}), tree importances, elastic net}\\[4pt]
  {\color{greenhl!45!black}\textbf{+} nearly free; interactions handled}\\[2pt]
  {\color{accent}\textbf{--} tied to one model family}};

\node[align=left, text width=12.8cm, font=\tiny, text=gray!55!black] at (4.5,-4.0)
     {\textbf{Practical recipe:} filter out the obviously useless (zero variance, $>90\%$
      missing, near-duplicates), then let an embedded method do the real work. Reserve
      wrappers for small feature sets where the extra compute is affordable. And run
      selection \emph{inside} the cross-validation folds --- selecting on the full dataset
      is leakage (\chref{wrangling}).};
\end{tikzpicture}}%
{Three families of feature selection. The choice is mostly a compute-versus-quality
trade-off, but the leakage warning applies to all three equally.}{feature-selection}

\section{Ensembles: Many Weak Models Beat One Strong One}

\begin{definitionbox}[Ensemble]
A model made of many models. Predictions are combined --- by voting for
classification, averaging for regression. Ensembles work because individual models
make \emph{different} mistakes, and averaging cancels the uncorrelated part of the
error.
\end{definitionbox}

\dsfig{%
\begin{tikzpicture}[font=\scriptsize,
  d/.style={rounded corners=2pt, draw=#1, fill=#1!10, text=#1!50!black,
            minimum width=1.5cm, minimum height=0.62cm, font=\tiny\bfseries},
  arr/.style={-{Stealth[length=1.8mm]}, gray!60, line width=0.7pt}]

% ---------------- BAGGING ----------------
\node[font=\scriptsize\bfseries, text=greenhl!45!black, anchor=west] at (-0.6,2.6)
     {BAGGING (e.g.\ Random Forest) --- models trained \emph{in parallel} on bootstrap samples};
\node[d=primary, minimum width=1.9cm] (data1) at (0.4,1.4) {full data};
\foreach \i in {1,2,3,4}{
  \pgfmathsetmacro{\xx}{2.6+(\i-1)*1.75}
  \node[d=greenhl] (s\i) at (\xx,1.9) {sample \i};
  \node[d=greenhl] (t\i) at (\xx,0.95) {tree \i};
  \draw[arr] (data1) -- (s\i);
  \draw[arr] (s\i) -- (t\i);}
\node[d=purplehl, minimum width=2.0cm] (vote) at (11.4,1.4) {vote / average};
\foreach \i in {1,2,3,4}{\draw[arr] (t\i) -- (vote);}
\node[font=\tiny\itshape, text=gray!55!black, text width=4.0cm, align=left] at (2.9,0.15)
     {each tree sees a different random sample \emph{and} a random subset of features
      $\Rightarrow$ their errors are decorrelated};
\node[font=\tiny\bfseries, text=greenhl!45!black, anchor=west] at (7.6,0.15)
     {reduces VARIANCE (fixes overfitting)};

\draw[gray!35, line width=0.5pt] (-0.7,-0.65) -- (12.7,-0.65);

% ---------------- BOOSTING ----------------
\node[font=\scriptsize\bfseries, text=accent!75!black, anchor=west] at (-0.6,-1.1)
     {BOOSTING (e.g.\ Gradient Boosting) --- models trained \emph{in sequence}, each fixing the last};
\node[d=primary, minimum width=1.9cm] (data2) at (0.4,-2.3) {full data};
\foreach \i in {1,2,3,4}{
  \pgfmathsetmacro{\xx}{2.9+(\i-1)*2.0}
  \node[d=accent] (b\i) at (\xx,-2.3) {tree \i};}
\draw[arr] (data2) -- (b1);
\foreach \i/\j in {1/2,2/3,3/4}{
  \draw[-{Stealth[length=1.8mm]}, accent!70, line width=0.9pt] (b\i) -- (b\j);}
\node[d=purplehl, minimum width=2.0cm] (sum) at (11.4,-2.3) {weighted sum};
\draw[arr] (b4) -- (sum);
\foreach \i in {1,2,3}{
  \pgfmathsetmacro{\xx}{2.9+(\i-1)*2.0}
  \node[font=\tiny, text=accent!70!black] at (\xx+1.0,-1.85) {residuals};}
\node[font=\tiny\itshape, text=gray!55!black, text width=5.2cm, align=left] at (3.5,-3.35)
     {tree $k+1$ is trained on what tree $k$ got \emph{wrong}, so the sequence
      cannot be parallelised};
\node[font=\tiny\bfseries, text=accent!75!black, anchor=west] at (9.0,-3.35)
     {reduces BIAS (fixes underfitting)};
\end{tikzpicture}}%
{Bagging versus boosting. Bagging builds independent models and averages away their
variance; boosting builds dependent models, each correcting its predecessor, and drives
down bias. This is the distinction to remember.}{bagging-boosting}

\subsection{Random Forests}

\begin{definitionbox}[Random Forest]
Bagging applied to decision trees, with one addition: at each split, only a random
subset of features may be considered. Without that, every tree would pick the same
dominant feature first and the trees would be near-identical --- defeating the
whole purpose.
\end{definitionbox}

\subsection{Gradient Boosting}

\begin{definitionbox}[Gradient Boosting]
Fit a weak model. Compute its residuals. Fit the next weak model to those
residuals. Add it to the ensemble, scaled by a \term{learning rate}. Repeat.
Implementations you will meet: XGBoost, LightGBM, CatBoost.
\end{definitionbox}

\rowcolors{2}{white}{defbg}
\begin{longtable}{@{}p{3.3cm}p{5.3cm}p{5.9cm}@{}}
\toprule
\rowcolor{primary!15}
& \textbf{Random Forest} & \textbf{Gradient Boosting} \\
\midrule
\endhead
Training & Parallel --- fast on many cores & Sequential --- slower \\
Tuning & Forgiving; defaults usually fine & Sensitive; needs learning rate, depth, $n$ trees \\
Overfitting & Resistant; more trees never hurts & Will overfit if over-trained; needs early stopping \\
Typical accuracy & Very good & Usually the best on tabular data \\
Choose it when & You want a strong result with little effort & You need maximum accuracy and can afford tuning \\
\bottomrule
\end{longtable}

\begin{conceptbox}[The Honest Default for Tabular Data]
Start with a random forest. It is hard to misuse, needs no scaling, handles mixed
types, and gives a strong number in one line. If that number is close enough to
your success criterion, stop. Move to gradient boosting only when the remaining gap
matters and you have the time to tune it properly.
\end{conceptbox}

\section{Reading Feature Importances --- Carefully}

\begin{alertbox}[Three Traps in Feature Importance]
\begin{tight}
  \item \textbf{Correlated features split their importance.} If attendance and
        lectures-watched correlate at 0.94, each may show half the importance it
        deserves, making both look unimportant.
  \item \textbf{Impurity-based importance is biased} towards high-cardinality
        features. A near-unique ID column can appear highly ``important''.
  \item \textbf{Importance is not causation.} A feature can be important to the
        model because it is a proxy for something you must not act on
        (\chref{ethics}).
\end{tight}
\textbf{Prefer permutation importance:} shuffle one column and measure how much
performance drops. It is slower, model-agnostic, and measures what you actually
mean by ``does this feature matter''.
\end{alertbox}

%---------------------------------------------------------------------
\begin{fourlenses}{Engineering Features}
\lensLA{The highest-value features are almost never raw counts. Build
\texttt{days\_since\_last\_activity}, the \emph{slope} of weekly engagement, and
\texttt{consecutive\_missed\_labs}. A student's trajectory carries far more signal
than their total. Use a random forest for accuracy, but keep a depth-3 tree beside
it as the explanation you can defend to an appeals panel.}

\lensPM{With 40 sprints, ensembles will overfit spectacularly. Engineer two or three
strong features instead --- \texttt{points\_committed / historical\_velocity},
\texttt{share\_of\_tasks\_with\_unresolved\_dependencies},
\texttt{team\_changed\_since\_last\_sprint} --- and fit something simple. This is a
domain where feature engineering does not merely beat algorithm choice; it is the
only thing that works.}

\lensIOT{Raw sensor values are weak; \emph{window statistics} are strong. Compute
rolling mean, standard deviation, min, max, and rate of change over 1-minute,
1-hour and 24-hour windows, plus frequency-domain features for vibration. This
single step usually matters more than anything in \chref{architectures}. Gradient
boosting on well-engineered window features is the standard, hard-to-beat baseline
for predictive maintenance.}

\lensSEC{Aggregate over time and entity: failed logins per account in 5 minutes,
distinct destination ports per source in 1 minute, bytes-out to bytes-in ratio,
deviation from that account's own 30-day baseline. The last is the key idea ---
features relative to an entity's \emph{own} history detect compromise far better
than absolute thresholds. Tree ensembles dominate here, but keep in mind that an
adversary who learns which features you use will manipulate them
(\chref{cybersecurity}).}
\end{fourlenses}

%---------------------------------------------------------------------
\begin{lab}[Engineer, Ensemble, Interpret]
\begin{lstlisting}[language=Python]
import pandas as pd, numpy as np
from sklearn.ensemble import RandomForestClassifier, HistGradientBoostingClassifier
from sklearn.inspection import permutation_importance
from sklearn.model_selection import cross_val_score

# --- 1. engineer from raw click logs -----------------------------------
g = clicks.groupby("student")
feats = pd.DataFrame({
    "n_clicks":      g.size(),
    "n_resources":   g.resource.nunique(),
    "n_active_days": g.timestamp.apply(lambda s: s.dt.date.nunique()),
    "days_since":    (CUTOFF - g.timestamp.max()).dt.days,      # recency: often the best
    "night_share":   g.timestamp.apply(lambda s: (s.dt.hour < 5).mean()),
    "daily_sd":      g.timestamp.apply(lambda s: s.dt.date.value_counts().std()),
})
# trend: is engagement rising or falling?
weekly = clicks.groupby(["student", clicks.timestamp.dt.isocalendar().week]).size()
feats["trend"] = weekly.groupby("student").apply(
    lambda s: np.polyfit(range(len(s)), s.values, 1)[0] if len(s) > 1 else 0.0)

# --- 2. ensembles ------------------------------------------------------
rf  = RandomForestClassifier(n_estimators=400, random_state=0, n_jobs=-1)
gb  = HistGradientBoostingClassifier(random_state=0)
for name, m in [("random forest", rf), ("grad boosting", gb)]:
    s = cross_val_score(m, X_tr, y_tr, cv=5, scoring="average_precision")
    print(f"{name:14s} PR-AUC {s.mean():.3f} +/- {s.std():.3f}")

# --- 3. interpret: permutation, NOT the built-in impurity importance ----
rf.fit(X_tr, y_tr)
perm = permutation_importance(rf, X_te, y_te, n_repeats=20,
                              random_state=0, scoring="average_precision")
order = perm.importances_mean.argsort()[::-1]
for i in order[:8]:
    print(f"  {feature_names[i]:24s} {perm.importances_mean[i]:+.4f}"
          f" +/- {perm.importances_std[i]:.4f}")
\end{lstlisting}
\end{lab}

%---------------------------------------------------------------------
\begin{checkpoint}
\begin{tightnum}
  \item Explain in one sentence each how bagging and boosting reduce error, and
        which component of error each targets.
  \item Why does a random forest sample features at each split as well as sampling
        rows?
  \item A tree ensemble reports \texttt{student\_id} as the third most important
        feature. What has gone wrong?
\end{tightnum}
\end{checkpoint}

%---------------------------------------------------------------------
\begin{chaptersummary}
\begin{tight}
  \item Good features beat clever algorithms. Ratios, deltas, window aggregations,
        recency and trend usually outperform raw columns.
  \item Feature selection comes in filter, wrapper and embedded flavours; run it
        inside the cross-validation folds or it leaks.
  \item Ensembles work because member models make different mistakes.
  \item Bagging (random forest) trains in parallel and reduces variance; boosting
        trains sequentially on residuals and reduces bias.
  \item Random forest is the forgiving default; gradient boosting usually wins on
        tabular data but needs tuning and early stopping.
  \item Use permutation importance, and remember importance is neither causation
        nor a licence to act.
\end{tight}
\end{chaptersummary}

\begin{reviewq}
  \item \textbf{(MCQ)} Boosting principally reduces: (a)~variance (b)~bias
        (c)~leakage (d)~dimensionality.
  \item \textbf{(MCQ)} Which requires no feature scaling? (a)~SVM (b)~k-NN
        (c)~random forest (d)~ridge regression.
  \item \textbf{(Short)} Explain the curse of dimensionality and its effect on
        distance-based methods.
  \item \textbf{(Short)} Why can two features that are individually useless be
        jointly decisive, and which selection family would miss this?
  \item \textbf{(Short)} State two reasons to prefer permutation importance over
        impurity-based importance.
  \item \textbf{(Applied)} You have one row per network connection with columns
        \texttt{src\_ip}, \texttt{dst\_port}, \texttt{bytes\_in},
        \texttt{bytes\_out}, \texttt{timestamp}. Engineer eight features for
        detecting compromised hosts, and say what each is meant to capture.
  \item \textbf{(Applied)} A team's random forest scores PR-AUC 0.71 and their
        gradient boosting machine 0.74 after two weeks of tuning. Advise them,
        referring to effort, maintenance and the success criterion.
\end{reviewq}
